\documentclass[11pt,a4paper]{article}

% Paquetes necesarios
\usepackage[utf8]{inputenc}
\usepackage[spanish]{babel}
\usepackage[margin=2cm]{geometry}
\usepackage{graphicx}
\usepackage{xcolor}
\usepackage{tcolorbox}
\usepackage{fontawesome5}
\usepackage{hyperref}
\usepackage{tikz}
\usetikzlibrary{shapes,arrows,positioning,shadows}
\usepackage{enumitem}
\usepackage{float}
\usepackage{fancyhdr}
\usepackage{titlesec}
\usepackage{multirow}

% Definir colores de Charlaton
\definecolor{charlatonBlue}{HTML}{0D5E9E}
\definecolor{charlatonCyan}{HTML}{24C4E8}
\definecolor{charlatonOrange}{HTML}{F7941D}
\definecolor{charlatonBg}{HTML}{F0F8FF}
\definecolor{charlatonGray}{HTML}{5A5A5A}

% Configurar hipervínculos
\hypersetup{
    colorlinks=true,
    linkcolor=charlatonBlue,
    urlcolor=charlatonCyan,
    pdftitle={Manual de Usuario - Charlaton},
    pdfauthor={Charlaton Team}
}

% Configurar encabezado y pie de página
\pagestyle{fancy}
\fancyhf{}
\fancyhead[L]{\color{charlatonBlue}\textbf{Charlaton}}
\fancyhead[R]{\color{charlatonGray}\textit{Manual de Usuario}}
\fancyfoot[C]{\color{charlatonGray}\thepage}
\renewcommand{\headrulewidth}{2pt}
\renewcommand{\headrule}{\hbox to\headwidth{\color{charlatonCyan}\leaders\hrule height \headrulewidth\hfill}}

% Estilo de secciones
\titleformat{\section}
{\color{charlatonBlue}\Large\bfseries}
{\thesection}{1em}{}[\color{charlatonCyan}\titlerule]

\titleformat{\subsection}
{\color{charlatonBlue}\large\bfseries}
{\thesubsection}{1em}{}

% Caja personalizada para consejos
\newtcolorbox{tipbox}{
    colback=charlatonBg,
    colframe=charlatonCyan,
    boxrule=1.5pt,
    arc=4pt,
    left=10pt,
    right=10pt,
    top=10pt,
    bottom=10pt
}

% Caja para advertencias
\newtcolorbox{warningbox}{
    colback=yellow!10,
    colframe=charlatonOrange,
    boxrule=1.5pt,
    arc=4pt,
    left=10pt,
    right=10pt,
    top=10pt,
    bottom=10pt
}

\begin{document}

% Portada
\begin{titlepage}
    \centering
    \vspace*{2cm}
    
    % Logo simulado con TikZ
    \begin{tikzpicture}
        \shade[ball color=charlatonBlue] (0,0) circle (1.5cm);
        \node[white,font=\Huge\bfseries] at (0,0) {C};
    \end{tikzpicture}
    
    \vspace{1cm}
    
    {\Huge\bfseries\color{charlatonBlue} Charlaton\par}
    \vspace{0.5cm}
    {\Large\color{charlatonGray} Plataforma de Videoconferencias\par}
    \vspace{2cm}
    
    {\huge\color{charlatonBlue} Manual de Usuario\par}
    \vspace{1cm}
    
    {\Large\color{charlatonGray} Versión 1.0\par}
    \vspace{0.5cm}
    {\large\color{charlatonGray} Diciembre 2024\par}
    
    \vfill
    
    % Iconos de características
    \begin{tikzpicture}
        \node[color=charlatonCyan] at (0,0) {\faVideo};
        \node[color=charlatonBlue] at (2,0) {\faComments};
        \node[color=charlatonOrange] at (4,0) {\faUsers};
        \node[color=charlatonCyan] at (6,0) {\faLock};
    \end{tikzpicture}
    
    \vspace{1cm}
    {\large\color{charlatonGray} Conéctate, comunica y colabora\par}
\end{titlepage}

% Tabla de contenidos
\tableofcontents
\newpage

% ==================== SECCIÓN 1: INTRODUCCIÓN ====================
\section{¡Bienvenido a Charlaton! \faSmile}

\begin{center}
\begin{tikzpicture}
    \node[draw=charlatonCyan,fill=charlatonBg,rounded corners=10pt,minimum width=12cm,minimum height=1.5cm,drop shadow] {
        \parbox{11cm}{\centering\large\color{charlatonBlue}
        \textbf{Charlaton} es tu plataforma de videoconferencias\\
        fácil, rápida y accesible para todos}
    };
\end{tikzpicture}
\end{center}

\vspace{0.5cm}

Charlaton te permite realizar reuniones virtuales con audio, video y chat en tiempo real. Ya sea para estudiar, trabajar, o simplemente charlar con amigos, ¡aquí lo puedes hacer todo!

\subsection{¿Qué puedes hacer en Charlaton?}

\begin{itemize}[leftmargin=2cm]
    \item[\color{charlatonCyan}\faVideo] \textbf{Videoconferencias:} Reuniones con audio y video de alta calidad
    \item[\color{charlatonBlue}\faComments] \textbf{Chat en tiempo real:} Mensajes instantáneos durante las reuniones
    \item[\color{charlatonOrange}\faUsers] \textbf{Salas compartidas:} Crea o únete a salas públicas o privadas
    \item[\color{charlatonCyan}\faLock] \textbf{Seguridad:} Salas protegidas con contraseña
    \item[\color{charlatonBlue}\faRecordVinyl] \textbf{Grabaciones:} Guarda tus reuniones importantes
\end{itemize}

\begin{tipbox}
    \faLightbulb\ \textbf{Consejo:} Charlaton funciona en cualquier navegador moderno. No necesitas instalar nada.
\end{tipbox}

% ==================== SECCIÓN 2: COMENZANDO ====================
\section{Primeros Pasos \faRocket}

\subsection{Crear tu cuenta}

\begin{enumerate}[leftmargin=2cm]
    \item \textbf{Ingresa a la plataforma:} Abre tu navegador y ve a la página de Charlaton
    \item \textbf{Regístrate:} Haz clic en \textcolor{charlatonBlue}{\textbf{"Registrarse"}}
    \item \textbf{Completa tus datos:}
    \begin{itemize}
        \item Correo electrónico
        \item Contraseña segura
        \item Nombre de usuario (nickname)
    \end{itemize}
    \item \textbf{¡Listo!} Ya puedes ingresar a la plataforma
\end{enumerate}

\begin{center}
\begin{tikzpicture}[node distance=1.5cm]
    \node[draw,rectangle,rounded corners,fill=charlatonBg,minimum width=3cm,minimum height=1cm,drop shadow] (step1) {\small\textbf{Registrarse}};
    \node[draw,rectangle,rounded corners,fill=charlatonBg,minimum width=3cm,minimum height=1cm,right=of step1,drop shadow] (step2) {\small\textbf{Completar datos}};
    \node[draw,rectangle,rounded corners,fill=charlatonBg,minimum width=3cm,minimum height=1cm,right=of step2,drop shadow] (step3) {\small\textbf{Comenzar}};
    \draw[thick,color=charlatonCyan,-latex] (step1) -- (step2);
    \draw[thick,color=charlatonCyan,-latex] (step2) -- (step3);
\end{tikzpicture}
\end{center}

\subsection{Iniciar Sesión}

Si ya tienes cuenta, simplemente:
\begin{enumerate}[leftmargin=2cm]
    \item Haz clic en \textcolor{charlatonBlue}{\textbf{"Iniciar Sesión"}}
    \item Ingresa tu correo y contraseña
    \item ¡Entra y disfruta!
\end{enumerate}

\begin{warningbox}
    \faExclamationTriangle\ \textbf{Importante:} Guarda tu contraseña en un lugar seguro. Si la olvidas, puedes recuperarla usando tu correo electrónico.
\end{warningbox}

% ==================== SECCIÓN 3: SALAS ====================
\section{Salas de Reunión \faDoorOpen}

\subsection{Crear una nueva sala}

\begin{enumerate}[leftmargin=2cm]
    \item En el menú principal, haz clic en \textcolor{charlatonBlue}{\textbf{"Crear Sala"}}
    \item Completa la información:
    \begin{itemize}
        \item \textbf{Nombre:} Ponle un nombre descriptivo a tu sala
        \item \textbf{Tipo:} Elige si será pública o privada
        \item \textbf{Contraseña:} (Opcional) Para salas privadas
        \item \textbf{Programar:} (Opcional) Agenda para más tarde
    \end{itemize}
    \item Haz clic en \textcolor{charlatonBlue}{\textbf{"Crear"}}
\end{enumerate}

\begin{tipbox}
    \faLightbulb\ \textbf{Consejo:} Las salas públicas aparecen en el listado general. Las privadas solo son accesibles con la contraseña.
\end{tipbox}

\subsection{Unirse a una sala existente}

\begin{center}
\begin{tikzpicture}
    % Opción 1
    \node[draw=charlatonBlue,fill=charlatonBg,rectangle,rounded corners,minimum width=5cm,minimum height=2.5cm,drop shadow] at (0,0) {
        \parbox{4.5cm}{
            \centering\color{charlatonBlue}\textbf{\faList\ Método 1}\\[0.3cm]
            \small\color{charlatonGray}
            1. Ver lista de salas\\
            2. Buscar la sala\\
            3. Clic en "Unirse"
        }
    };
    % Opción 2
    \node[draw=charlatonOrange,fill=yellow!10,rectangle,rounded corners,minimum width=5cm,minimum height=2.5cm,drop shadow] at (6,0) {
        \parbox{4.5cm}{
            \centering\color{charlatonOrange}\textbf{\faLink\ Método 2}\\[0.3cm]
            \small\color{charlatonGray}
            1. Obtén el enlace\\
            2. Ábrelo en navegador\\
            3. Ingresa contraseña si aplica
        }
    };
\end{tikzpicture}
\end{center}

\vspace{0.5cm}

\textbf{Salas con contraseña:} Si la sala es privada, se te pedirá la contraseña antes de entrar.

% ==================== SECCIÓN 4: EN LA REUNIÓN ====================
\section{Durante la Reunión \faVideo}

\subsection{Panel de Control}

Una vez dentro de la sala, verás el panel de control con varios botones:

\begin{center}
\begin{tabular}{cl}
\textcolor{charlatonCyan}{\faMicrophone} & \textbf{Micrófono:} Activar/Desactivar audio \\[0.3cm]
\textcolor{charlatonBlue}{\faVideo} & \textbf{Cámara:} Activar/Desactivar video \\[0.3cm]
\textcolor{charlatonOrange}{\faComments} & \textbf{Chat:} Abrir/Cerrar el chat \\[0.3cm]
\textcolor{charlatonCyan}{\faUsers} & \textbf{Participantes:} Ver lista de usuarios \\[0.3cm]
\textcolor{red}{\faPhoneSlash} & \textbf{Salir:} Abandonar la reunión \\
\end{tabular}
\end{center}

\begin{tipbox}
    \faLightbulb\ \textbf{Consejo:} Presiona los botones para activar o desactivar. Los iconos cambiarán de color para indicar el estado.
\end{tipbox}

\subsection{Video y Audio}

\textbf{Para activar tu cámara o micrófono:}
\begin{itemize}
    \item Haz clic en el botón correspondiente
    \item La primera vez, tu navegador te pedirá permiso
    \item Autoriza el acceso a la cámara/micrófono
\end{itemize}

\begin{warningbox}
    \faExclamationTriangle\ \textbf{Importante:} Si no ves video o no escuchas audio, verifica que tu navegador tenga permisos para acceder a la cámara y micrófono.
\end{warningbox}

\textbf{Controlar tu audio/video:}
\begin{itemize}
    \item \textcolor{charlatonCyan}{\faMicrophone}: Micrófono activado
    \item \textcolor{red}{\faMicrophoneSlash}: Micrófono silenciado
    \item \textcolor{charlatonBlue}{\faVideo}: Cámara activada
    \item \textcolor{red}{\faVideoSlash}: Cámara desactivada
\end{itemize}

\subsection{Chat en Tiempo Real}

El chat te permite enviar mensajes a todos los participantes:

\begin{enumerate}[leftmargin=2cm]
    \item Haz clic en el botón \textcolor{charlatonOrange}{\faComments}
    \item Escribe tu mensaje en el campo de texto
    \item Presiona Enter o el botón de enviar
    \item ¡Tu mensaje aparecerá para todos!
\end{enumerate}

\begin{tipbox}
    \faLightbulb\ \textbf{Consejo:} El chat permanece visible durante toda la reunión y puedes revisar mensajes anteriores haciendo scroll.
\end{tipbox}

\subsection{Participantes}

Para ver quién está en la sala:
\begin{itemize}
    \item Haz clic en \textcolor{charlatonCyan}{\faUsers}
    \item Verás la lista completa de participantes
    \item Los iconos indican si tienen audio/video activo
\end{itemize}

% ==================== SECCIÓN 5: FUNCIONES AVANZADAS ====================
\section{Funciones Adicionales \faStar}

\subsection{Administración de Salas}

Si eres el \textbf{creador de la sala}, tienes privilegios especiales:

\begin{itemize}[leftmargin=2cm]
    \item[\color{charlatonBlue}\faCrown] \textbf{Finalizar reunión:} Cierra la sala para todos
    \item[\color{charlatonCyan}\faUserPlus] \textbf{Agregar administradores:} Comparte el control
    \item[\color{charlatonOrange}\faKey] \textbf{Cambiar contraseña:} Actualiza la seguridad
    \item[\color{charlatonBlue}\faTrash] \textbf{Eliminar sala:} Borra la sala permanentemente
\end{itemize}

\subsection{Grabación de Reuniones}

\begin{warningbox}
    \faExclamationTriangle\ \textbf{Nota:} La función de grabación está disponible para administradores de la sala.
\end{warningbox}

Para grabar una reunión:
\begin{enumerate}[leftmargin=2cm]
    \item Haz clic en el botón de grabar \textcolor{red}{\faCircle}
    \item La grabación comenzará automáticamente
    \item Para detener, presiona el mismo botón
    \item La grabación se guardará en tu cuenta
\end{enumerate}

\subsection{Compartir Sala}

Para invitar a otros:
\begin{enumerate}[leftmargin=2cm]
    \item Copia el enlace de la sala desde tu navegador
    \item Envíalo por correo, WhatsApp, etc.
    \item Si la sala tiene contraseña, compártela también
\end{enumerate}

\begin{center}
\begin{tikzpicture}
    \node[draw=charlatonCyan,fill=charlatonBg,rounded corners,minimum width=10cm,minimum height=1cm,drop shadow] {
        \parbox{9cm}{\centering\small\texttt{https://charlaton.com/room/abc123xyz}}
    };
\end{tikzpicture}
\end{center}

% ==================== SECCIÓN 6: CONSEJOS ====================
\section{Consejos para una Mejor Experiencia \faThumbsUp}

\subsection{Conexión y Calidad}

\begin{itemize}[leftmargin=2cm]
    \item[\color{charlatonCyan}\faWifi] \textbf{Internet estable:} Usa WiFi o conexión por cable
    \item[\color{charlatonBlue}\faDesktop] \textbf{Cierra otras aplicaciones:} Libera recursos del navegador
    \item[\color{charlatonOrange}\faVolumeUp] \textbf{Usa audífonos:} Evita eco y mejora el audio
\end{itemize}

\subsection{Privacidad y Seguridad}

\begin{warningbox}
    \faLock\ \textbf{Seguridad:} Nunca compartas tus contraseñas. Cierra sesión en computadoras compartidas.
\end{warningbox}

\begin{itemize}[leftmargin=2cm]
    \item Usa contraseñas únicas y seguras
    \item No compartas enlaces de salas privadas públicamente
    \item Reporta comportamientos inapropiados
\end{itemize}

\subsection{Etiqueta en Reuniones}

\begin{center}
\begin{tikzpicture}
    \node[draw=charlatonBlue,circle,fill=charlatonBg,minimum size=2cm,drop shadow] at (0,0) {\Large\color{charlatonBlue}\faMicrophoneSlash};
    \node[below] at (0,-1.5) {\small\parbox{2.5cm}{\centering\textbf{Silencia cuando no hables}}};
    
    \node[draw=charlatonCyan,circle,fill=charlatonBg,minimum size=2cm,drop shadow] at (4,0) {\Large\color{charlatonCyan}\faComments};
    \node[below] at (4,-1.5) {\small\parbox{2.5cm}{\centering\textbf{Usa el chat sin interrumpir}}};
    
    \node[draw=charlatonOrange,circle,fill=yellow!10,minimum size=2cm,drop shadow] at (8,0) {\Large\color{charlatonOrange}\faSmile};
    \node[below] at (8,-1.5) {\small\parbox{2.5cm}{\centering\textbf{Sé respetuoso siempre}}};
\end{tikzpicture}
\end{center}

% ==================== SECCIÓN 7: RESOLUCIÓN DE PROBLEMAS ====================
\section{Solución de Problemas \faTools}

\subsection{No puedo activar mi cámara o micrófono}

\textbf{Solución:}
\begin{enumerate}[leftmargin=2cm]
    \item Verifica los permisos del navegador
    \item Ve a Configuración → Privacidad y seguridad
    \item Asegúrate de permitir acceso a cámara/micrófono
    \item Recarga la página
\end{enumerate}

\subsection{No veo a otros participantes}

\textbf{Solución:}
\begin{itemize}
    \item Verifica tu conexión a internet
    \item Recarga la página
    \item Intenta salir y volver a entrar a la sala
\end{itemize}

\subsection{El chat no funciona}

\textbf{Solución:}
\begin{itemize}
    \item Verifica que estás conectado a la sala
    \item Revisa tu conexión a internet
    \item Actualiza la página
\end{itemize}

\begin{tipbox}
    \faQuestionCircle\ \textbf{¿Necesitas más ayuda?} Visita nuestra sección de soporte o contacta al equipo de Charlaton.
\end{tipbox}

% ==================== SECCIÓN 8: PREGUNTAS FRECUENTES ====================
\section{Preguntas Frecuentes \faQuestion}

\subsection{¿Cuántas personas pueden estar en una sala?}

Charlaton soporta reuniones con múltiples participantes. El número exacto depende de tu plan y conexión.

\subsection{¿Necesito descargar algo?}

No. Charlaton funciona completamente en tu navegador. Solo necesitas un navegador moderno como Chrome, Firefox o Edge.

\subsection{¿Es gratis?}

Charlaton ofrece un plan gratuito con funcionalidades básicas. También hay planes premium con funciones avanzadas.

\subsection{¿Funciona en móvil?}

Sí. Puedes usar Charlaton desde tu teléfono o tablet abriendo la página en tu navegador móvil.

\subsection{¿Las reuniones son seguras?}

Sí. Todas las comunicaciones están cifradas y puedes proteger tus salas con contraseña.

% ==================== SECCIÓN 9: CONTACTO ====================
\section{Contacto y Soporte \faEnvelope}

\begin{center}
\begin{tcolorbox}[colback=charlatonBg,colframe=charlatonBlue,width=12cm,arc=10pt,boxrule=2pt]
    \centering
    \Large\color{charlatonBlue}\textbf{¿Necesitas ayuda?}\\[0.5cm]
    \normalsize\color{charlatonGray}
    Visita nuestra página de soporte\\[0.2cm]
    o contáctanos por correo:\\[0.3cm]
    \textcolor{charlatonCyan}{\texttt{soporte@charlaton.com}}\\[0.5cm]
    \small Estamos aquí para ayudarte 24/7
\end{tcolorbox}
\end{center}

% ==================== CONCLUSIÓN ====================
\section*{¡Comienza Ahora! \faRocket}
\addcontentsline{toc}{section}{¡Comienza Ahora!}

\begin{center}
\begin{tikzpicture}
    \node[draw=charlatonBlue,fill=charlatonBg,rounded corners=15pt,minimum width=13cm,minimum height=3cm,drop shadow,line width=2pt] {
        \parbox{12cm}{
            \centering
            \Large\color{charlatonBlue}\textbf{¡Ya estás listo para usar Charlaton!}\\[0.5cm]
            \normalsize\color{charlatonGray}
            Crea tu cuenta, únete a salas o crea las tuyas propias.\\
            Conéctate con amigos, familia, compañeros de estudio o trabajo.\\[0.3cm]
            \large\color{charlatonCyan}\textbf{¡La comunicación empieza aquí!}
        }
    };
\end{tikzpicture}
\end{center}

\vspace{1cm}

\begin{center}
\begin{tikzpicture}
    \node[color=charlatonCyan,scale=2] at (0,0) {\faVideo};
    \node[color=charlatonBlue,scale=2] at (2,0) {\faComments};
    \node[color=charlatonOrange,scale=2] at (4,0) {\faUsers};
    \node[color=charlatonCyan,scale=2] at (6,0) {\faHeart};
\end{tikzpicture}
\end{center}

\vfill

\begin{center}
\color{charlatonGray}\small
---\\
Manual de Usuario - Charlaton v1.0\\
© 2024 Charlaton Team. Todos los derechos reservados.
\end{center}

\end{document}
